% !PRINTED 180
\weberpage{180}
Bemerkenswert sind auch die Differentialgleichungen:
\weq{9}{d\varkappa^2 = -d\varkappa'^2 = \pi i\,\vartheta_{00}^4\,\varkappa^2\varkappa'^2\,d\omega\ [\S\,52,\,(3)],}
\weq{10}{d\varkappa^2\varkappa'^2 = -\pi i\,\vartheta_{00}^4(\varkappa^2-\varkappa'^2)\varkappa^2\varkappa'^2\,d\omega,}
\weq{11}{df(\omega)^8 = \dfrac{\pi i}{3}\vartheta_{00}^4[f_2(\omega)^8 - f_1(\omega)^8]\,d\omega,}
\weq{12}{d\gamma_2(\omega) = -\dfrac{2\pi i}{3}\gamma_3(\omega)\eta(\omega)^4\,d\omega\ [\text{nach}\ (5)\ \text{und}\ \S\,34,\,(10)],}
\weq{13}{dj(\omega) = -2\pi i\,\gamma_2(\omega)^2\gamma_3(\omega)\eta(\omega)^4\,d\omega.}
Man erhält sodann für die linearen Fundamentaltransformationen aus \S\,34, (13), (14), (15):
\weq{14}{\begin{array}{ll}
\gamma_2(\omega+1) = e^{-\frac{2\pi i}{3}}\gamma_2(\omega), & \gamma_2\left(-\dfrac{1}{\omega}\right) = \gamma_2(\omega),\\[6pt]
\gamma_3(\omega+1) = -\gamma_3(\omega), & \gamma_3\left(-\dfrac{1}{\omega}\right) = -\gamma_3(\omega),
\end{array}}
und aus \S\,40 findet man allgemein:
\weq{15}{\begin{array}{l}
\gamma_2\left(\dfrac{\gamma+\delta\omega}{\alpha+\beta\omega}\right) = \varrho\,\gamma_2(\omega),\\[6pt]
\gamma_3\left(\dfrac{\gamma+\delta\omega}{\alpha+\beta\omega}\right) = (-1)^{\alpha\gamma+\beta\gamma+\beta\delta}\gamma_3(\omega),
\end{array}}
worin $\varrho$ wie in \S\,40, (3) die Bedeutung hat:
\[
\varrho = e^{-\frac{2\pi i}{3}(\alpha\gamma+\beta\delta-\alpha\beta-\alpha^2\beta\delta)}.
\]
Es sind dies also Modulfunktionen, und die Gruppen, zu denen sie gehören, sind durch die beiden Kongruenzen charakterisiert:
\[
\begin{array}{ll}
-\alpha\beta+\alpha\gamma+\beta\delta-\alpha^2\beta\gamma &\equiv 0\ (\mathrm{mod}\ 3),\\
\alpha\gamma+\beta\gamma+\beta\delta &\equiv 0\ (\mathrm{mod}\ 2).
\end{array}
\]
Wir geben hiernach unserem Satz 4. die folgende Ergänzung, die sich aus den Formeln (14) ergibt.

5. \gesperrt{Eine Modulfunktion, die durch die beiden Substitutionen}
\[
(\omega, \omega+1), \qquad \left(\omega, -\dfrac{1}{\omega}\right)
\]
\gesperrt{das Zeichen ändert, ist das Produkt von $\gamma_3(\omega)$ mit einer rationalen Funktion von $j(\omega)$.}
